\subsection{将函数解释为缓增广义函数}

定义 5. --- 设 \(\nu\) 是 \(\mathbf{R}^{n}\) 上的测度。如果存在整数 \(r \in \mathbf{N}\) ，使得连续映射 \(x \mapsto (1 + \|x\|)^{-r}\) 关于 \(\nu\) 在 \(\mathbf{R}^{n}\) 上可积，就称它是缓增的。

换言之，测度 \(\nu\) 是缓增的，如果存在 \(r \in \mathbf{N}\) ，使得由 \(x \mapsto \| x\|^{-r}\) 所定义的函数关于 \(\nu\) 在 \(\mathbf{R}^n\) 中单位球的余集上可积。特别地， \(\mathbf{R}^n\) 上的每个有界测度都是缓增的。更一般地，若 \(f\) 是多项式增长的 \(\mu\)-可测函数，且 \(\nu\) 是缓增的，则测度 \(f \cdot \nu\) 是缓增的。

\(\mathcal{M}^{\mathrm{t}}(\mathbf{R}^{n})\) （\(\mathbf{R}^{n}\) 上缓增测度所成的集）是 \(\mathcal{M}(\mathbf{R}^{n};\mathbf{C})\) （\(\mathbf{R}^{n}\) 上复测度所成的空间）的一个向量子空间。

命题 15. --- 设 \(\nu\) 是 \(\mathbf{R}^{n}\) 上的缓增测度。 \(\nu\) 在 \(\mathcal{S}(\mathbf{R}^{n})\) 上的限制是一个缓增广义函数。它为零当且仅当测度 \(\nu\) 为零。

由于 \(\nu\) 是缓增的，存在正整数 \(k\) ，使得映射 \(x \mapsto \| x\|^{-k}\) 关于 \(\nu\) 在 \(\mathbf{R}^n\) 中单位球的余集上可积。对每个施瓦茨函数 \(\varphi \in \mathcal{S}(\mathbf{R}^n)\) ，我们有

\[
| \langle \nu , \varphi \rangle | \leqslant \Bigl (\int_ {\| x \| \leqslant 1} d \nu \Bigr) \widetilde {q} _ {0, 0} (\varphi) + \Bigl (\int_ {\| x \| > 1} \| x \| ^ {- k} d \nu \Bigr) \widetilde {q} _ {0, k} (\varphi),
\]

因此 \(\nu\) 在 \(\mathcal{S}(\mathbf{R}^n)\) 上的限制是一个缓增广义函数。

最后的断言由 IV, p. 206 的命题 6 得出，因为 \(\mathcal{D}(\mathbf{R}^n)\) 包含在 \(\mathcal{S}(\mathbf{R}^n)\) 中。

我们将把空间 \(\mathcal{M}^{\mathrm{t}}(\mathbf{R}^{n})\) 等同于 \(\mathcal{S}'(\mathbf{R}^{n})\) 的一个子空间。

命题 16. --- 设 \(p \in [1, +\infty]\) ， \(f \in \mathcal{L}^{p}(\mathbf{R}^{n})\) 。那么测度 \(f \cdot \mu\) （以 \(f\) 为密度，关于勒贝格测度）是缓增的。映射 \(f \mapsto f \cdot \mu\) （从 \(\mathrm{L}^{p}(\mathbf{R}^{n})\) 到 \(\mathcal{S}'(\mathbf{R}^{n})\) 中的）是连续的。

设 \(q\) 是 \(p\) 的共轭指数，并设 \(r \geqslant 0\) 满足 \(rq > n\) 。对每个 \(x \in \mathbf{R}^n\) ，令 \(g(x) = (1 + \|x\|)^{-r}\) 。函数 \(g\) 属于 \(\mathcal{L}^q(\mathbf{R}^n)\) （由 IV, p. 199 的命题 3）。由赫尔德不等式，我们有

\[
\int_ {\mathbf {R} ^ {n}} ^ {*} (1 + \| x \|) ^ {- r} | f (x) | d \mu (x) \leqslant \mathrm{N} _ {q} (g) \mathrm{N} _ {p} (f) <   + \infty ,
\]

因此测度 \(f \cdot \mu\) 是缓增的。

设 \(f \in \mathrm{L}^p(\mathbf{R}^n)\) ， \(\varphi \in \mathcal{S}(\mathbf{R}^n)\) 。赫尔德不等式蕴含

\[
| \langle f \cdot \mu , \varphi \rangle | = \left| \int_ {\mathbf {R} ^ {n}} f (x) \varphi (x) d \mu (x) \right| \leqslant \| f \| _ {p} \| \varphi \| _ {q}
\]

而映射 \(f \mapsto f \cdot \mu\) 的连续性则由 IV, p. 213 的命题 13 得出。

对每个 \(p \in [1, +\infty]\) ，我们将把 \(\mathrm{L}^p(\mathbf{R}^n)\) 等同于 \(\mathcal{S}'(\mathbf{R}^n)\) 的一个子空间（借助线性映射 \(f \mapsto f \cdot \mu\) ）。

命题 17. --- 空间 \(\mathcal{D}(\mathbf{R}^{n})\) 与 \(\mathcal{S}(\mathbf{R}^{n})\) 在 \(\mathcal{S}'(\mathbf{R}^{n})\) 中稠密。

只须证明 \(\mathcal{D}(\mathbf{R}^{n})\) 在 \(\mathcal{S}'(\mathbf{R}^{n})\) 中稠密。设 \(\lambda\) 是 \(\mathcal{S}'(\mathbf{R}^{n})\) 上的连续线性形式，它在 \(\mathcal{D}(\mathbf{R}^{n})\) 上为零。由于空间 \(\mathcal{S}(\mathbf{R}^{n})\) 是自反的，存在函数 \(\varphi \in \mathcal{S}(\mathbf{R}^{n})\) ，使得对所有 \(\lambda(f) = \langle f, \varphi \rangle\) 有 \(f \in \mathcal{S}'(\mathbf{R}^{n})\) 。于是我们有

\[
0 = \lambda (\psi) = \langle \psi , \varphi \rangle = \int_ {\mathbf {R} ^ {n}} \psi \varphi d \mu
\]

对所有 \(\psi\in\mathcal{D}(\mathbf{R}^{n})\) 成立。因此 \(\varphi\cdot\mu\) 上的测度 \(\mathbf{R}^{n}\) 为零（IV, p. 206 的命题 6），从而 \(\varphi=0\) 。命题于是由哈恩–巴拿赫定理（EVT, II, p. 46，推论 1）得出。

